\chapter{Commutation Failure}
\label{chap:commutation-failure}

% Closed question: What records the failure of transport to commute?

\begin{chapteressay}

Some operations commute. Adding three to four is the same as adding
four to three. Multiplying two by five is the same as multiplying
five by two. The order does not matter; the result is determined by
the inputs alone.

Other operations do not commute. Rotating a die ninety degrees about
the $x$-axis and then ninety degrees about the $y$-axis produces a
different orientation than the reverse order. Applying Alice's code
change and then Bob's code change produces a different final program
than the reverse order. Running constant folding before dead code
elimination produces different optimized code than the reverse
order. In each case, the algorithm's result depends on the sequence,
not just the operations.

This chapter is about the difference. The difference is a structured
quantity called the commutator. The commutator captures, in an
algebraic form, the failure of operations to commute. Two
operations $A$ and $B$ have commutator $[A, B] = AB - BA$. If $A$
and $B$ commute, the commutator is zero. If they do not commute, the
commutator is nonzero, and its value carries information about how
the order matters.

The chapter develops five algorithmic forms of the commutator. The
first is the matrix commutator: rotations in three-dimensional space
have a non-zero commutator algebra (the Lie algebra so(3)). The
second is the git merge conflict: two developers making different
changes to the same code produce a commutation failure that the
version-control tool flags as a conflict. The third is compiler
optimization order: LLVM's optimization passes have specified
commutation properties; the commutator records the order dependence.
The fourth is the Jacobi identity: an algebraic constraint on the
commutator's behavior on triple products, ensuring the algebra
closes consistently. The fifth is \texttt{HeartbeatProcess} in
Lean: the formal record of how proof elaboration cost depends on
the order of inference rule applications.

The second concrete example is the print shop. A four-color print
job applies cyan, magenta, yellow, and black inks in sequence. The
order of application matters: yellow over black produces a
different appearance than black over yellow. The press operator's
craft is managing the order to achieve the desired result. The
non-commutation of color applications is a real, observable
property of the print. The print operator is constantly working
with commutators, even though the term is not used in the trade.

The chapter's closed question is:

\begin{center}
\emph{What records the failure of transport to commute?}
\end{center}

The answer is the commutator. The commutator is a derived
quantity: it is computed from the two operations and their
composition. It has the same operational type as the original
operations: it can be applied to records, it can be composed with
other operations, it has its own commutators (the Jacobi identity
constrains how these closed-form commutators relate).

The chapter's algorithmic claim is that the commutator is the
admissible record of non-commutation. When the algorithm encounters
two operations that produce different results in different orders,
the algorithm does not arbitrarily pick one order; it records the
commutator. The record is part of the algorithm's output: the
result of $AB$, the result of $BA$, and their difference.

In some contexts (rotation groups, gauge theories, neural network
training), the commutator is a structured algebraic object with
well-known properties. In other contexts (git merges, compiler
optimizations, print shop sequences), the commutator is a less
formal but equally real quantity: the visible difference between
two orderings.

The chapter's contribution is to identify this universal pattern.
Any algorithm that operates on order-sensitive inputs has
commutators. The commutators are part of the algorithm's record.
The algorithm cannot avoid them; it can only record them and
manage them.

Chapter 10 takes up the question of invariance: what does survive
admissible order rearrangement? The commutators tell us what
changes; the invariants tell us what does not. Together, the two
chapters bracket the algorithmic structure of admissible
computation.

\end{chapteressay}

\section{The Matrix Commutator}
\label{se:the-matrix-commutator}

Rotations in three-dimensional space do not commute. Rotating an object
ninety degrees about the $x$-axis and then ninety degrees about the
$y$-axis produces a different orientation than rotating ninety degrees
about the $y$-axis and then ninety degrees about the $x$-axis. The
difference is observable: pick up a die, rotate it as described, and
note the face on top. The two orderings produce different top faces.

In matrix terms, let $R_x$ be the rotation matrix for ninety degrees
about the $x$-axis and $R_y$ be the rotation matrix for ninety degrees
about the $y$-axis. The compositions $R_x R_y$ and $R_y R_x$ are
different matrices. Their difference is the commutator:
\[
[R_x, R_y] = R_x R_y - R_y R_x.
\]
For the specific case of ninety-degree rotations, the commutator is a
particular rotation matrix; for infinitesimal rotations, the
commutator is approximately $\theta^2 R_z$ where $\theta$ is the
small angle and $R_z$ is the rotation about the $z$-axis. The
commutator of two rotations is a third rotation, in the direction
perpendicular to both.

This is the chapter's question in its simplest form. Two operations
that should be admissible separately produce different results when
composed in different orders. The difference is the commutator. The
commutator is itself an operation: a rotation about a specific axis.
The non-commutation is a structural fact about rotations.

The matrix commutator is the algebraic core of the algorithmic
chapter. Wherever two operations do not commute, the difference is
captured by their commutator. The commutator measures the failure
of order-independence. If the operations commuted, the commutator
would be zero; the algorithm could perform them in any order. If
they do not commute, the commutator is nonzero; the algorithm must
choose an order, and the choice is part of the algorithm's output.

For three-dimensional rotations, the commutator algebra is so(3), a
three-dimensional Lie algebra. The commutators are the structural
constants: $[J_x, J_y] = J_z$, $[J_y, J_z] = J_x$, $[J_z, J_x] = J_y$.
The algebra closes: every commutator is itself in the algebra. The
closure is the algebraic admissibility: the chapter's structure
remains within the same operational space.

In Lean, the matrix commutator appears in the
\texttt{CalculusProcess} typeclass in \texttt{Episode11.lean}. The
class formalizes the operation of taking derivatives along paths.
Two derivatives in different directions do not commute in general;
their commutator is the chapter's algorithmic record of the
non-commutation. The class carries the commutator as a derived
quantity.

The chapter's algorithmic claim is that any algorithm that applies
two operations to the same data, and whose result depends on the
order of application, has a non-zero commutator. The commutator is
the data structure that records the dependence on order. The
record is admissible: it is a derived value of the same type as
the original operations. It enters the algorithm's output as a
structured residue.

\section{Git Merge Conflict}
\label{se:git-merge-conflict}

Two developers work on the same software project. Both clone the
repository, both create branches from the same commit, both edit the
same file. Alice modifies line 40, changing the variable name from
\texttt{x} to \texttt{count}. Bob modifies the same line 40,
changing the type annotation from \texttt{int} to \texttt{float}.
Each commits their change to their own branch.

Now the team merges. The merge tool fetches Alice's branch and Bob's
branch and tries to combine them. For most lines in the file, the two
branches are identical: they were branched from the same commit and
modified independently. For line 40, the two branches disagree.
Neither change can be silently dropped. The merge tool reports a
conflict.

This is a commutation failure. Alice's change followed by Bob's
change does not produce the same result as Bob's change followed by
Alice's change. The forward direction of Alice's change (renaming
\texttt{x} to \texttt{count}) and the forward direction of Bob's
change (changing the type) operate on different aspects of the same
line, but the resulting line differs based on order: if Alice goes
first, the new line is ``count: int = 0''; Bob then changes the
type to give ``count: float = 0''. If Bob goes first, the new line
is ``x: float = 0''; Alice then renames to give ``count: float = 0''.

Wait --- in this case, both orderings produce ``count: float = 0''.
The changes do commute. The merge tool, however, does not have a
semantic model of the code. It sees two textual changes to the same
line and reports a conflict because it cannot decide which textual
form is correct without human intervention.

Let me change the example. Alice modifies line 40 to insert a new
statement: \texttt{x = x + 1}. Bob modifies line 40 to insert a
different statement: \texttt{x = x * 2}. After Alice's change, the
sequence is ``$x = x + 1$''. After Bob's change applied to that,
the sequence becomes ``$x = x + 1; x = x * 2$''. The resulting value
of $x$ is $(x + 1) \cdot 2 = 2x + 2$.

If Bob's change goes first, the sequence is ``$x = x * 2$''. After
Alice's change applied to that, the sequence is ``$x = x * 2; x = x
+ 1$''. The resulting value of $x$ is $2x + 1$.

The two orderings produce different final values. The two changes
do not commute. The merge tool reports a conflict because the
human must decide which order produces the desired behavior. The
conflict resolution is the human's choice of order; the resolution
becomes the canonical merged record.

This is the chapter's structure in version control. Two operations
(Alice's change, Bob's change) applied to the same record (the
line of code) produce different results in different orders. The
difference is the commutator. The merge tool's conflict report is
the commutator: it identifies which lines have non-commuting
changes and marks them for human resolution.

In Lean, the analog appears in the \texttt{SpaceTimePath.le}
relation in \texttt{Episode11.lean}. The relation defines a partial
order on paths through the structural space. Two paths from the
same origin to the same destination may or may not commute. When
they do not commute, the commutator records the difference. The
order is the merge resolution: choose one path as canonical, and
treat the other as a deviation requiring justification.

The git merge conflict is the chapter's everyday example. Software
development teams encounter it constantly. The resolution requires
human judgment: which change is correct, which should be retained,
which represents the team's intention. The tool cannot decide; the
team must.

This is a foreshadowing of the chapter's final algorithmic
observation. The commutator measures the failure of commutativity.
When two algorithms produce different outputs in different orders,
the difference is not arbitrary; it is structured. The structure
is the commutator's specific value. The algorithm can compute the
commutator; the algorithm cannot decide which value is correct.
The decision requires an external agent (the team, the developer,
the user) to choose between the two orders.

\section{Compiler Optimization Order}
\label{se:compiler-optimization-order}

The LLVM compiler infrastructure applies a sequence of optimization
passes to compiled code. Each pass is a transformation: dead code
elimination removes unreachable statements; constant folding
replaces expressions like \texttt{2 + 3} with \texttt{5}; loop
unrolling converts \texttt{for (i=0; i<3; i++) f(i)} into
\texttt{f(0); f(1); f(2)}; inlining replaces function calls with
the function body.

The order in which the passes are applied matters. Two passes may
produce different final code depending on which runs first. The
LLVM manual documents this carefully: certain passes are required
to run before others; some pairs are recommended in a specific
order; some pairs commute and can be run in either order. The
commutation is part of the optimization pipeline's specification.

Consider constant folding and dead code elimination. Constant
folding replaces \texttt{2 + 3} with \texttt{5}. Dead code
elimination removes \texttt{if false then x = 5}. The two passes
do not commute: constant folding might convert \texttt{if (2 +
3 > 4) then x = 5} into \texttt{if true then x = 5}, which is
admissible code, not dead. Dead code elimination first would not
recognize the conditional as dead because the condition contains
unfolded constants. Constant folding first exposes the
condition as always true, allowing the dead-code pass to simplify
the \texttt{if true} structure.

The non-commutation is a structural fact about the optimization
algebra. The compiler's pipeline is therefore order-sensitive.
Different pipelines produce different optimized code. The same
source, compiled with two different optimization orders, may
result in two different binaries.

The optimization commutator is the difference between the two
results. If we denote constant folding as $C$ and dead code
elimination as $D$, the commutator is $[C, D] = CD - DC$. For
most code, the commutator is small (the two orders produce
nearly identical results); for some code, the commutator is
significant (the two orders produce different control flow,
different register allocation, different cache behavior).

This is the chapter's algorithmic measure in compiler optimization.
The compiler's pipeline is a specific sequence of operations. The
sequence is the algorithm's commitment to one ordering. The
commutator of the sequence measures the deviation from the
opposite ordering. The compiler's manual specifies which orderings
are admissible (which produce semantically equivalent code) and
which are not (which may produce different output).

In Lean, the analog appears in the \texttt{YarnTheory.le} relation
in \texttt{Episode11.lean}. The relation defines a partial order
on operations in the yarn structure (the chapter's term for the
weave of paths). Two operations in different parts of the yarn may
or may not commute. The relation records which pairs commute and
which do not. The commutator of non-commuting pairs is the
chapter's structural quantity.

The compiler analogy makes the chapter's structure concrete. The
algorithm's order of operations is part of its specification.
Different orders may produce different results. The commutator
measures the difference. The compiler's documentation is the
algorithm's commitment to specific orders for specific reasons.

This connects to broader algorithmic discipline. Functional
programming languages encourage referential transparency: the order
of operations should not matter. Pure functions commute trivially.
Imperative programming, by contrast, is order-dependent: statements
have side effects that depend on their order of execution. The
chapter's claim is that order-dependence is admissible and
algorithmically tractable, but it must be recorded explicitly. The
commutator is the record.

\section{Jacobi Identity}
\label{se:jacobi-identity}

The Jacobi identity is an algebraic constraint on the commutator
operation in a Lie algebra. For any three elements $X$, $Y$, $Z$ of
the algebra:
\[
[[X, Y], Z] + [[Y, Z], X] + [[Z, X], Y] = 0.
\]
The identity says that the commutator's commutators close in a
specific way: the cyclic sum is zero. The Jacobi identity is what
distinguishes a Lie algebra from an arbitrary algebra with a
bilinear bracket. Without it, the commutator's structure would not
close consistently.

For the Lie algebra sl(2) of $2 \times 2$ traceless complex
matrices, the generators are:
\[
E = \begin{pmatrix} 0 & 1 \\ 0 & 0 \end{pmatrix}, \quad
F = \begin{pmatrix} 0 & 0 \\ 1 & 0 \end{pmatrix}, \quad
H = \begin{pmatrix} 1 & 0 \\ 0 & -1 \end{pmatrix}.
\]
The commutators are $[E, F] = H$, $[H, E] = 2E$, $[H, F] = -2F$.
The Jacobi identity can be verified by computing each of the
three triple commutators and summing.

The verification proceeds case by case. Take $X = E$, $Y = F$, $Z =
H$. Then:
\begin{align*}
[[E, F], H] &= [H, H] = 0, \\
[[F, H], E] &= [-2F \cdot (-1), E] = [2F, E] = -2 [E, F] = -2H, \\
[[H, E], F] &= [2E, F] = 2[E, F] = 2H.
\end{align*}
The sum is $0 + (-2H) + 2H = 0$. The identity holds.

This is the algorithmic check the chapter requires of any
commutator algebra. The commutator captures the non-commutation;
the Jacobi identity verifies that the algebra of commutators
closes consistently. Without the Jacobi identity, the algebra
would have inconsistent triple products: the triple commutator
of three elements would not equal the cyclic sum, and the algebra
would not have well-defined associativity.

The Jacobi identity is the algorithmic version of consistency
verification. For a commutator algebra to be admissible as a
structural framework, it must satisfy the identity. Any algebra
that violates the identity is internally inconsistent: the
commutator's behavior on triple products is ambiguous.

The compiler's analog is the verification of pass commutation
properties. Before a compiler is shipped, its developers verify
that the optimization passes commute or non-commute as expected.
Any pass that violates the documented commutation properties is
a bug. The Jacobi-like verification (consistency of pass
interactions) is part of the compiler's quality assurance.

In a broader sense, the Jacobi identity is the algebraic
consistency condition for any operation that takes pairs of
records and produces a third record. The condition ensures that
the operation's structure is consistent across triples. The
condition is automatic for many natural algebraic structures
(Lie algebras, Lie groups, gauge theories); it is non-trivial for
ad-hoc operation collections, and its verification is the
algebraic check that the collection is well-formed.

In Lean, the Jacobi identity for \texttt{YarnTheory} operations
is a proof obligation. Any new operation introduced into the
yarn algebra must be verified against the existing operations
to satisfy the identity. The compiler checks the proof at
elaboration; a violation produces an error. The discipline is
strict: the algebra must close.

\section{HeartbeatProcess in Lean}
\label{se:heartbeatprocess-in-lean}

\texttt{Episode11.lean} defines a typeclass \texttt{HeartbeatProcess}
that records how elaboration cost accumulates along a proof path. The
typeclass is:

\begin{leanexcerpt}{HeartbeatProcess}
structure HeartbeatProcess
    (Value : Type)
    (Carrier : CarrierProcess Value)
    [DISTINGUISHABLE Value Carrier] ... [UNIVERSAL Value Carrier]
  where
  bullshit_meter       : CalculusProcess Value Carrier
  current_reading      : SpaceTimePath
  accumulated_bullshit : YarnTheory
  weave? : YarnTheory $\to$ YarnTheory := ...
\end{leanexcerpt}

A \texttt{HeartbeatProcess} carries the full project cascade
through \texttt{UNIVERSAL} plus three fields: the
\texttt{bullshit\_meter} (the underlying \texttt{CalculusProcess}
that supplies the heartbeat reference), the
\texttt{current\_reading} along the spacetime path, and the
\texttt{accumulated\_bullshit} encoded as a \texttt{YarnTheory}
record. The \texttt{weave?} function advances the accumulated
yarn by one layer. As the proof elaborates, the structure updates
the yarn rather than a flat integer.

The interesting fact is that the total heartbeat count for a proof
can depend on the order in which the steps are taken. Two equivalent
proof strategies --- the same set of inference rules applied to reach
the same conclusion --- may have different total heartbeat counts if
they apply the rules in different orders. The difference is the
commutator of the rules' costs.

Consider a small example. Suppose the proof requires applying
\texttt{simp} and \texttt{rfl}. If \texttt{simp} is applied first, it
performs a complex term reduction that costs 100 heartbeats; the
subsequent \texttt{rfl} is cheap (10 heartbeats), checking the
already-simplified term. Total: 110 heartbeats. If \texttt{rfl} is
applied first, it costs 50 heartbeats (checking the original
unreduced term); \texttt{simp} on the post-\texttt{rfl} state may
either succeed cheaply or fail entirely if \texttt{rfl} produced an
incompatible representation. Total: 100 (succeed) or infinity (fail).

The two strategies are not commutatively equivalent. The order
matters. The \texttt{HeartbeatProcess} records the order's effect:
the accumulated cost differs by strategy. The compiler does not
prefer one strategy over the other automatically; the user must
choose.

This is the chapter's structure in formal Lean tooling. Two
admissible paths to the same conclusion may have different costs.
The cost is the chapter's algorithmic distance from start to
conclusion. The non-commutation of the rules means the distance
depends on the path. The commutator is the difference.

The \texttt{HeartbeatProcess} is the algorithmic record of the
non-commutation. The compiler's elaboration tracks the heartbeat
count and decides whether the elaboration is admissible at the
current rate (using the \texttt{maxHeartbeats} option). If the
elaboration exceeds the budget, the compiler stops and reports a
failure. The order-dependence of the heartbeat count means that
slight changes in the proof strategy can move the elaboration
from below-budget to above-budget. The cost is order-dependent;
the admissibility decision is order-dependent.

The \texttt{YarnTheory.le} relation extends this to a partial order
on the yarn-theoretic representations of paths. The chapter's term
``yarn'' captures the woven structure of multiple paths: each
strand is a proof path, the strands cross at shared lemmas, the
weave is the algebraic structure of the proof system. The order
relation says: one yarn weave is at most another if it can be
extended to the other by adding strands without disturbing the
existing weave.

The commutator in this context is the failure of two strands to
combine into a single thread. If two paths in the yarn cross
without commuting, the crossing is a knot. The knot is the
commutator's algebraic record. The yarn theory carries the
knots as named structural elements.

The chapter's algorithmic claim is that this is the Lean
formalization of the curvature recording. The
\texttt{HeartbeatProcess} records the cost of paths. The
\texttt{YarnTheory.le} records the structural relation among paths.
The combination records the commutators of the paths. The
compiler enforces the discipline: every elaboration that produces
a non-zero commutator records that commutator in the elaboration
log. The log is the chapter's algorithmic data structure for
non-commutation.

\begin{bridge}

The chapter's question was what records the failure of transport to
commute.

The answer is the commutator. Two transports applied in different
orders may produce different results; the commutator is the algebraic
difference. The matrix commutator captures rotational non-commutation.
Git merge conflicts record textual non-commutation. Compiler
optimization passes have specified commutation properties; the
commutator records when the order matters. The Jacobi identity
verifies that the commutator algebra closes consistently. The
\texttt{HeartbeatProcess} in Lean records the cost-dependence on
the elaboration path; \texttt{YarnTheory.le} records the structural
non-commutation in proof yarn.

What remains is invariance. Some quantities survive any admissible
relabeling of the computation. These are the invariants. Chapter 10
identifies the invariants of admissible computation: what does not
change under recompilation, gauge transformation, or relabeling?

\end{bridge}

\begin{coda}{Print-Shop Color Layers}

The press operator has been at the shop for thirty-one years. He knows
the machines, the inks, the papers, the press rooms, and the
peculiarities of each color job. Today he is running a four-color
print on coated stock: a poster for a museum exhibition, two thousand
copies, twenty-four by thirty-six inches.

Four-color process printing uses cyan, magenta, yellow, and black inks
(CMYK). Each color is printed by a separate plate; the press passes
the paper through four plate stations in sequence, laying down one
color at a time. The final image is the combination of the four
layers.

The order of the four colors is not arbitrary. Different orders
produce different visible results. The standard order in offset
printing is CMYK: cyan, then magenta, then yellow, then black. The
order is chosen because of the inks' physical properties: yellow is
the most transparent, so it sits on top of the other colors. Black is
last because it carries the image's dark detail and benefits from
being applied as the top opaque layer.

If the press is reconfigured to print in a different order --- say,
KCMY --- the visible result differs. The colors mix differently. The
yellow that should sit on top is now underneath the black, where it
contributes less brightness. The dark areas are slightly different in
hue. The image's overall tone shifts.

This is a commutation failure in the most literal sense. The
operations (applying each color) are admissible separately. The
composition depends on order. The commutator is the difference
between the two prints.

For the museum poster, the standard order has been used: CMYK. The
press operator has set up the plates accordingly. The first run is
made. The press operator inspects the print.

The first print has a problem. The black has a slight greenish cast
in the dark areas. The cause is a non-commutation between the yellow
and black: the yellow ink had not dried completely before the black
plate was applied, and the wet yellow mixed slightly with the black,
producing a green tint where the two overlapped.

The fix is to slow the press, allowing more drying time between plate
stations. The press operator reduces the press speed from 8,000
sheets per hour to 6,500 sheets per hour. The reduced speed gives
the yellow ink more time to dry before the black is applied. The
second run is made. The print's dark areas now have the correct
neutral black.

The change in press speed is a recalibration of the ordering. The
press is still applying CMYK, but the temporal spacing between the
applications is longer. The commutator between adjacent colors is
smaller (less wet-on-wet interaction). The print's tonal balance
improves.

This is the chapter's structure in printing form. The order of
operations matters; the commutator records the failure of
commutation; recalibration adjusts the operations' temporal spacing
to reduce the commutator. The print operator's craft is a continuous
management of these commutators: which color goes first, how much
drying time, how much pressure, how much ink density.

A more dramatic example. The museum poster's design includes a
varnish layer applied after the four-color print. Varnish is
transparent; it adds gloss and protection. The varnish can be
applied as a full coverage (uniform gloss) or as a spot varnish
(gloss only on specific areas, dull on the rest).

The order of the varnish matters. If the varnish is applied before
the CMYK process, the inks have a different surface to adhere to;
they spread differently, dry differently, and produce different
colors. If the varnish is applied after CMYK, the inks have already
set and the varnish seals them in. Standard practice is varnish
after CMYK.

But the museum has requested a specialty effect: the poster's title
should be printed in metallic gold ink, applied as a fifth color
plate. The gold ink is opaque and has its own drying properties.
Where in the sequence should it go?

The press operator considers. Gold ink is heavy and has a relatively
high coverage. If applied first, the subsequent CMYK plates will
deposit on top of it, but the gold's metallic shimmer may be muted
by the overlying colors. If applied last (after CMYK, before varnish),
the gold sits on top of the CMYK and has full visibility, but its
relatively high coverage may not adhere well to the dried CMYK.

The decision: gold ink is applied after black (after the standard
CMYK) but before varnish. The sequence becomes CMYK-gold-varnish:
six plate stations. The press is reconfigured. The plates are
loaded.

The press is now performing a five-color print with one varnish
finish. Each color is one admissible step. The sequence is six
operations. The commutator between any pair captures the order
dependence: gold-before-cyan would produce a different visual
than gold-after-cyan; varnish-before-gold would alter the gold's
finish.

Run the print. Inspect. The gold's metallic effect is preserved.
The varnish gives uniform gloss across the entire print. The
poster looks correct.

The next batch of two hundred fifty prints is run. The press is
slowed slightly more to accommodate the extra plate stations.
Each print emerges as a faithful copy of the design.

This is the print operator's daily work: managing the commutators
of a multi-stage process. Each stage is a transformation of the
paper. Each combination of stages has its specific commutation
properties. The operator selects the order that produces the
desired result, knowing that other orders would produce different
results, knowing the difference is the commutator's specific
value.

The exhibition opens in three weeks. The poster will hang in the
museum's entrance hall for the duration of the show. The print's
quality will be inspected by museum staff, the artist, and the
public. The print operator's management of the commutators has
produced a result that all three constituencies will find
satisfactory. The commutators are real; their management is craft.

\end{coda}
